\chapter{Anhang}
\label{chap:Anhang}

\section{Glossar}

\begin{description}
\item[AES (Advanced Encryption Standard)] Symmetrisches Verschlüsselungsverfahren, das Daten blockweise in mehreren Runden verändert.
\item[ARPANET] Vorläufer des heutigen Internets, ein Forschungsnetzwerk der USA.
\item[ASCII] Zeichenkodierung, die Zahlen Buchstaben und Zeichen zuordnet.
\item[Authentifizierung] Prüfung, ob Daten tatsächlich vom angegebenen Absender stammen und nicht verändert wurden.
\item[Backend und Frontend] Backend bezeichnet den Teil einer Software, der im Hintergrund Daten verarbeitet. Das Frontend ist die Benutzeroberfläche.
\item[Bibliothek (Library)] Sammlung von fertigem Programmcode, die in eigenen Programmen verwendet werden kann.
\item[Brute Force] Angriffsmethode, bei der alle möglichen Schlüssel systematisch durchprobiert werden.
\item[Chiffretext] Verschlüsselter, ohne Schlüssel unlesbarer Text.
\item[Client] Programm, das eine Verbindung zu einem Server aufbaut und dessen Dienste nutzt.
\item[Client-Server-Architektur] Aufbau eines Netzwerks, bei dem Clients ihre Anfragen an einen zentralen Server senden.
\item[Combobox] Auswahlfeld einer GUI, in dem man einen Eintrag aus einer Liste wählt.
\item[Daemon Thread] Thread, der im Hintergrund läuft und automatisch endet, sobald alle normalen Threads beendet sind.
\item[Datenpaket] Abgeschlossene Dateneinheit, die über ein Netzwerk verschickt wird.
\item[Diffie-Hellman-Merkle-Schlüsselaustausch] Verfahren, mit dem zwei Teilnehmer über einen offenen Kanal einen gemeinsamen geheimen Schlüssel berechnen.
\item[DNS (Domain Name System)] Dienst, der Domainnamen in IP-Adressen umwandelt.
\item[EAX-Modus] Betriebsmodus von AES, der Daten verschlüsselt und gleichzeitig authentifiziert.
\item[Endlicher Körper (Galois-Körper)] Mathematische Struktur mit endlich vielen Elementen, in der man addieren, subtrahieren, multiplizieren und dividieren kann. AES rechnet in \(\mathbb{F}_{2^8}\) mit 256 Elementen.
\item[Exploit] Programm oder Methode, die eine Sicherheitslücke ausnutzt.
\item[Flag] Statusvariable mit den Werten wahr oder falsch.
\item[Fork] Abspaltung eines bestehenden Softwareprojekts, die eigenständig weiterentwickelt oder veröffentlicht wird.
\item[Frame] Datenpaket der Sicherungsschicht mit Header und Trailer.
\item[GUI (Graphical User Interface)] Grafische Benutzeroberfläche.
\item[Header und Trailer] Zusatzinformationen vor und nach den Nutzdaten eines Datenpakets, zum Beispiel Adressen.
\item[HTTPS] Protokoll der Anwendungsschicht zur verschlüsselten Übertragung von Webseiten.
\item[Hybride Verschlüsselung] Kombination aus asymmetrischem Schlüsselaustausch und anschliessender symmetrischer Verschlüsselung.
\item[IP (Internet Protocol)] Protokoll der Vermittlungsschicht, das Datenpakete mit IP-Adressen versieht und weiterleitet.
\item[Indikator] Erstes Byte des eigenen Protokolls, das den Pakettyp angibt (0x0 Liste, 0x1 Nachricht, 0x2 Fehlercode).
\item[IP-Adresse] Eindeutige Adresse eines Geräts in einem Netzwerk.
\item[Kryptogramm] Verschlüsselter Nachrichteninhalt im eigenen Protokoll.
\item[Kryptographie] Wissenschaft der Verschlüsselung von Informationen.
\item[Lock] Sperre, die verhindert, dass mehrere Threads gleichzeitig auf dieselben Daten zugreifen.
\item[MAC-Adresse (Media Access Control)] Eindeutige Hardwareadresse einer Netzwerkkarte.
\item[Man-in-the-Middle (MitM)] Angriff, bei dem sich ein Angreifer unbemerkt in die Kommunikation zweier Teilnehmer einschaltet.
\item[Matrix] Rechteckige Anordnung von Zahlen in Zeilen und Spalten.
\item[Multithreading] Gleichzeitige Ausführung mehrerer Threads in einem Programm.
\item[Nonce (number used once)] Zufällige Zahl, die bei demselben Schlüssel nie wiederverwendet werden darf.
\item[Nutzdaten (Payload)] Eigentlicher Dateninhalt eines Pakets ohne Header und Trailer.
\item[OSI-Modell] Referenzmodell, das die Netzwerkkommunikation in sieben Schichten gliedert.
\item[Port] Nummer, die auf einem Gerät einem bestimmten Programm oder Dienst zugeordnet ist.
\item[Protokoll] Vereinbarung über Aufbau und Ablauf eines Datenaustauschs.
\item[Puffer (Buffer)] Speicherbereich, der eingehende Daten vorübergehend aufnimmt.
\item[Race Condition] Fehler, der entsteht, wenn das Ergebnis eines Programms vom zeitlichen Ablauf mehrerer Threads abhängt.
\item[Router] Gerät, das Datenpakete zwischen Netzwerken weiterleitet.
\item[Rundenschlüssel] Teilschlüssel, der aus dem AES-Schlüssel abgeleitet und in einer Runde verwendet wird.
\item[S-Box (Substitutionsbox)] Tabelle, die in AES jedes Byte durch ein anderes ersetzt.
\item[Schlüsselexpansion] Berechnung der Rundenschlüssel aus dem ursprünglichen AES-Schlüssel.
\item[Server] Programm oder Gerät, das Dienste für Clients bereitstellt.
\item[Socket] Schnittstelle des Betriebssystems für die Kommunikation über ein Netzwerk.
\item[Switch und Bridge] Geräte der Sicherungsschicht, die Netzwerkteilnehmer anhand von MAC-Adressen verbinden.
\item[Symmetrische und asymmetrische Verschlüsselung] Bei der symmetrischen Verschlüsselung verwenden beide Seiten denselben Schlüssel, bei der asymmetrischen ein Schlüsselpaar aus öffentlichem und privatem Schlüssel.
\item[Tag] Prüfwert des EAX-Modus, mit dem der Empfänger kontrolliert, ob die Nachricht unverändert ist.
\item[TCP (Transmission Control Protocol)] Zuverlässiges Protokoll der Transportschicht mit Three-Way-Handshake, Sequenzierung und Checksumme.
\item[Thread] Ausführungsstrang eines Programms.
\item[Tkinter] Python-Bibliothek für grafische Benutzeroberflächen.
\item[UDP (User Datagram Protocol)] Schnelles, aber weniger zuverlässiges Protokoll der Transportschicht.
\item[UTF-8] Zeichenkodierung, die Zeichen als Bytefolgen darstellt.
\item[uv] Werkzeug zur Verwaltung von Python-Projekten und deren Abhängigkeiten.
\item[Widget] Einzelnes Bedienelement einer GUI, zum Beispiel Schaltfläche, Eingabefeld oder Beschriftung.
\item[Wireshark] Software, die Datenpakete eines Netzwerks mitschneidet und anzeigt.
\item[XOR (exklusives Oder)] Logische Verknüpfung, die genau dann wahr ist, wenn sich die beiden Eingangswerte unterscheiden.
\item[Zotero und BibTeX] Programme zur Verwaltung von Quellen und zur Erstellung des Literaturverzeichnisses.
\end{description}

\section{Aufbau des Protokolls}
Das eigene Protokoll liegt in den Nutzdaten der TCP-Verbindung. Das erste Byte, der Indikator, bestimmt den Pakettyp.
\begin{table}[H]
    \centering
    \begin{tabular}{|l|c|c|l|}\hline
         Pakettyp & Indikator & Aufbau & Richtung\\\hline
         Liste & 0x0 & Indikator, danach je 6 Byte pro Benutzer (4 Byte IP, 2 Byte Port) & Server $\rightarrow$ Client\\\hline
         Nachricht & 0x1 & Indikator, IP, Port, Tag, Nonce, Kryptogramm & Client $\rightarrow$ Server $\rightarrow$ Client\\\hline
         Fehlercode & 0x2 & Nur Indikator, Empfänger nicht gefunden & Server $\rightarrow$ Client\\\hline
    \end{tabular}
    \caption{Pakettypen des eigenen Protokolls}
    \label{tab:A1}
\end{table}
\begin{table}[H]
    \centering
    \begin{tabular}{|l|c|c|}\hline
         Feld & Länge in Byte & Position (ab 0 gezählt)\\\hline
         Indikator & 1 & 0\\\hline
         IP-Adresse & 4 & 1 bis 4\\\hline
         Port & 2 & 5 bis 6\\\hline
         Tag & 16 & 7 bis 22\\\hline
         Nonce & 16 & 23 bis 38\\\hline
         Kryptogramm & Länge der Nachricht & ab 39\\\hline
    \end{tabular}
    \caption{Felder eines Nachrichtenpakets (0x1)}
    \label{tab:A2}
\end{table}

\section{Quellcode}
Der vollständige Quellcode aller Versionen liegt im GitHub-Repository: \url{https://github.com/Luka-Adjancic/Luka-Adjancic-Mauturaarbeit}.